\section{Integral polynomial lattices and normalized Fock duality}
\label{coef:section}

We now construct the coefficient lattices used at the beginning of the
finite-group argument in Section~\ref{div:section}, and prove their
duality. The submitted construction uses polynomial dual-Weyl and Weyl
lattices, then transports them by local bases.\cite{alpo-uniform-source}
The proof below uses Gaussian elimination, divided row operations and
the elementary theory of lattices. In particular, the integral duality
will be proved here, rather than imported from a theorem on tableaux.

\subsection{The polynomial modules and their normalization}

Fix integers $0\le n\le r$ and a partition
$\nu=(\nu_1,\ldots,\nu_n)$, so
$\nu_1\ge\cdots\ge\nu_n\ge0$. Let $X=(x_{as})$ be an $r$ by $n$
matrix of independent variables. For a characteristic-zero domain $R$,
let $\nabla_\nu(R)$ consist of the polynomials in $R[X]$ which are
homogeneous of degree $\nu_s$ in the variables of each column $s$ and satisfy
\begin{equation}\label{coef:column-invariance}
 F\bigl(X(I+tE_{ij})\bigr)=F(X)\qquad(1\le i<j\le n)
\end{equation}
as polynomial identities in the additional indeterminate $t$.
Here $E_{ij}$ is the elementary matrix: the substitution adds $t$ times
column $i$ to column $j$. Identity \eqref{coef:column-invariance} is an
identity of polynomials, not merely a condition at the elements of a
residue field. Write $V_\nu=\nabla_\nu(\mathbb Q)$.

Let $\Delta_j(X)$ be the determinant of the first $j$ rows and first
$j$ columns, put $\nu_{n+1}=0$, and define
\[
 H_\nu(X)=\prod_{j=1}^n\Delta_j(X)^{\nu_j-\nu_{j+1}}.
\]
The polynomial $H_\nu$ belongs to $\nabla_\nu(\mathbb Z)$: every
initial minor is fixed by the indicated column additions, and the
degrees telescope to $\nu_s$. Its row degrees are
$(\nu_1,\ldots,\nu_n,0,\ldots,0)$.
Let $I_{r,n}$ denote the first $n$ columns of the identity matrix.
Then $H_\nu(I_{r,n})=1$.

For $1\le a,b\le r$ define the divided row operator
\[
 D_{ab}^{(k)}F=[t^k]F\bigl((I+tE_{ab})X\bigr),\qquad k\ge0.
\]
Thus the substitution changes row $a$ to row $a+t\,b$ when $a\ne b$.
For $a=b$ it multiplies row $a$ by $1+t$.
All these operators preserve integral coefficients. They preserve
$\nabla_\nu(R)$ because row substitutions commute with the right
column substitutions, and preserve column degrees. Let
\[
 \Delta_\nu(R)=\operatorname{span}_R
 \{D_1\cdots D_mH_\nu:m\ge0,\ D_i=D_{a_i b_i}^{(k_i)}\}.
\]
The empty product is allowed. In particular
$\Delta_\nu(R)\subset\nabla_\nu(R)$.
These are the polynomial Weyl and dual-Weyl lattices needed here.

For $F=\sum_\alpha F_\alpha X^\alpha$ and
$G=\sum_\alpha G_\alpha X^\alpha$ over $\mathbb Q$, define the
bilinear Fock pairing and its normalization by
\begin{equation}\label{coef:fock}
 [F,G]=\sum_\alpha\alpha!F_\alpha G_\alpha,
 \qquad C_\nu=[H_\nu,H_\nu],\qquad
 B_\nu(F,G)=\frac{[F,G]}{C_\nu},
 \qquad \alpha!=\prod_{a,s}\alpha_{as}!.
\end{equation}
The integer $C_\nu$ is positive: the Fock form is positive definite
on real polynomials, and $H_\nu(I_{r,n})=1$.
Thus the denominator in \eqref{coef:fock} is nonzero, including when
all parts of $\nu$ vanish. This explicit finite formula specifies the
normalizing scalar completely; a product formula for it is unnecessary.

\subsection{Integral coefficients from Gaussian elimination}

\begin{lemma}[Restriction to lower triangular matrices]\label{coef:restriction}
Let $F\in V_\nu$. Its restriction to matrices with $x_{as}=0$ for
$a<s$ determines $F$. If the restricted polynomial has integer
coefficients, then $F\in\mathbb Z[X]$.
\end{lemma}
\begin{proof}
The case $n=0$ concerns constant polynomials and is immediate. Assume
$n>0$ and put $\Delta=\prod_{j=1}^n\Delta_j$.
In the ring $A=\mathbb Z[X,\Delta^{-1}]$, Gaussian elimination gives
\[
                         X=LU,
\]
where $L$ is lower triangular in the rectangular sense
($L_{as}=0$ if $a<s$) and $U$ is upper triangular with diagonal entries
one. All entries belong to $A$.
For completeness, eliminate the entries above the diagonal in column
$j$, using the already treated columns $1,\ldots,j-1$ in this order.
The pivot in column $i$ is $\Delta_i/\Delta_{i-1}$, with
$\Delta_0=1$: these column operations leave each relevant initial
minor unchanged, and a lower triangular initial block has determinant
equal to the product of its pivots. Each pivot is a unit in $A$.
This proves the asserted factorization by induction, without dividing
by an integer.

An upper triangular unipotent matrix is a product of elementary
matrices $I+tE_{ij}$, $i<j$: successively clear its entries above the
diagonal by column additions. Therefore
\eqref{coef:column-invariance}, applied over the localized coefficient
ring, gives $F(X)=F(L)$. A zero restriction consequently implies $F=0$.
If the restriction has integer coefficients, it gives
$F\in A$.

We claim that $\mathbb Q[X]\cap A=\mathbb Z[X]$.
Write $F=G/m$ with $G\in\mathbb Z[X]$, $m>0$, and with no prime
dividing both $m$ and every coefficient of $G$. Membership in $A$
gives $\Delta^kG=mQ$ for some $k\ge0$ and $Q\in\mathbb Z[X]$.
For any prime $p\mid m$, the polynomial $\Delta$ is nonzero modulo
$p$: each initial determinant contains its diagonal monomial with
coefficient one, and a polynomial ring over $\mathbb F_p$ is a domain.
Reduction of the equality modulo $p$ therefore forces every coefficient
of $G$ to be divisible by $p$, a contradiction. Thus $m=1$.
\end{proof}

\begin{lemma}[The extremal row weight]\label{coef:extremal-weight}
The subspace of $V_\nu$ of row degree
$(\nu_1,\ldots,\nu_n,0,\ldots,0)$ is $\mathbb QH_\nu$.
Moreover, for every $F\in V_\nu$,
\begin{equation}\label{coef:evaluation-pairing}
                         B_\nu(H_\nu,F)=F(I_{r,n}).
\end{equation}
\end{lemma}
\begin{proof}
Projecting onto a fixed row degree preserves the column-invariance
identities, because every term of a column substitution preserves the
degree in each row. Let $F_0$ be the component of $F$ of the displayed
row degree.

A monomial with these row degrees and column degrees $\nu$, and
supported on entries with $a\ge s$, must be
\[
                         \prod_{s=1}^n x_{ss}^{\nu_s}.
\]
Indeed, row one can use only column one, so its prescribed degree
already exhausts that column; induction gives the assertion for every
row and column. The restriction of $H_\nu$ is the same diagonal
monomial. Lemma~\ref{coef:restriction} shows that
$F_0=cH_\nu$. Evaluation at $I_{r,n}$ leaves only the diagonal
monomial, so $c=F(I_{r,n})$.
Distinct row degrees are orthogonal for the Fock pairing, which is
diagonal on monomials. Since $H_\nu$ has the indicated row degree,
$[H_\nu,F]=[H_\nu,F_0]=cC_\nu$, proving
\eqref{coef:evaluation-pairing}.
\end{proof}

\begin{lemma}[Adjoint row operators]\label{coef:adjoint}
For a matrix $g\in\operatorname{Mat}_r(\mathbb Q)$, put
$T_gF(X)=F(gX)$. Then
\[
 [T_gF,G]=[F,T_{g^t}G],\qquad
 [D_{ab}^{(k)}F,G]=[F,D_{ba}^{(k)}G].
\]
The first identity holds over every characteristic-zero field after
extending the pairing, and applies independently to several disjoint
blocks of variables.
\end{lemma}
\begin{proof}
Constants satisfy the assertion directly. On polynomials of homogeneous
degree $m\ge1$ in any finite set of variables,
the multinomial theorem gives
\[
 [ (u\mathbin{\cdot}x)^m,(v\mathbin{\cdot}x)^m]
                         =m!(u\mathbin{\cdot}v)^m.
\]
The powers of linear forms span that space in characteristic zero.
One direct verification is the polarization identity
\[
 m!y_1\cdots y_m
  =\sum_{S\subset\{1,\ldots,m\}}(-1)^{m-|S|}
                               \left(\sum_{i\in S}y_i\right)^m:
\]
expanding the right side cancels every term missing one of the $m$
symbols and leaves the product term with coefficient $m!$.
Repetitions among the $y_i$ give every monomial. Substituting a matrix
in a linear form transposes its coefficient vector, so the first
identity follows from the displayed Fock formula and bilinearity.
Different total degrees are orthogonal. Apply this to the block matrix
which acts on the row variables of every column. Finally replace $g$
by $I+tE_{ab}$ and compare coefficients of $t^k$. This includes
$a=b$; in that case the divided operator multiplies a row-homogeneous
polynomial of row degree $d$ by $\binom dk$.
\end{proof}

For a word $D=D_1\cdots D_m$ let
$D^*=D_m^*\cdots D_1^*$, where
$(D_{ab}^{(k)})^*=D_{ba}^{(k)}$.
Lemmas~\ref{coef:extremal-weight} and \ref{coef:adjoint} give the
key identity
\begin{equation}\label{coef:word-evaluation}
                 B_\nu(F,DH_\nu)=(D^*F)(I_{r,n})
                       \qquad(F\in V_\nu).
\end{equation}

\subsection{Exact integral duality}

For a subgroup $L\subset V_\nu$ define
\[
 L^{\vee,B}=\{F\in V_\nu:B_\nu(F,L)\subset\mathbb Z\}.
\]

\begin{theorem}[Integral Weyl duality]\label{coef:integral-duality}
For every $0\le n\le r$ and every partition $\nu$ of length $n$,
both $\nabla_\nu(\mathbb Z)$ and $\Delta_\nu(\mathbb Z)$ are full
lattices in $V_\nu$, and
\begin{equation}\label{coef:duality}
 \Delta_\nu(\mathbb Z)^{\vee,B}=\nabla_\nu(\mathbb Z),
 \qquad
 \nabla_\nu(\mathbb Z)^{\vee,B}=\Delta_\nu(\mathbb Z).
\end{equation}
Thus $B_\nu$ gives a perfect integral pairing between these two
lattices. No restriction on the primes dividing $C_\nu$ is required.
\end{theorem}
\begin{proof}
Equation~\eqref{coef:word-evaluation} shows that the normalized pairing
of an integral polynomial in $\nabla_\nu$ with every generator of
$\Delta_\nu(\mathbb Z)$ is integral. This proves one inclusion in
the first equality.

Conversely, suppose $F\in V_\nu$ pairs integrally with
$\Delta_\nu(\mathbb Z)$. Then $(DF)(I_{r,n})\in\mathbb Z$ for
every word of divided row operators, since transposition and reversal
permute the set of all words. Consider the matrix $L(z)$ with diagonal
entries one, entries $z_{as}$ below the diagonal, and zero entries
above the diagonal. It is the first $n$ columns of a product of
elementary row matrices. More explicitly,
\[
 L(z)=\left(\prod_{s=1}^n\prod_{a>s}(I+z_{as}E_{as})\right)I_{r,n},
\]
where the factors are written from left to right in increasing $s$;
the order for fixed $s$ is immaterial. Products of two of the
off-diagonal summands vanish in this order: for a factor in column
$s$ followed by one in column $t\ge s$, a nonzero product would
require $s=b>t$. Consequently the displayed product has exactly the
specified entries. Expanding the substitutions one factor at a time
shows that every coefficient of $F(L(z))$ is $(DF)(I_{r,n})$ for an
appropriate word $D$, and hence is an integer.

Column homogeneity now recovers the restriction to an arbitrary lower
triangular matrix. Write its diagonal entries as $d_s$ and its lower
entries as $l_{as}$. In the Laurent polynomial ring in the $d_s$,
\[
 F((l_{as}))=\left(\prod_s d_s^{\nu_s}\right)
              F\bigl(L((l_{as}/d_s)_{a>s})\bigr).
\]
Every coefficient on the right is an integer. For each off-diagonal
monomial, the column degree determines its diagonal powers uniquely;
there is no collision of distinct terms. The left side is a polynomial,
so every surviving diagonal exponent is nonnegative. Its coefficients
are therefore integers. Lemma~\ref{coef:restriction} gives
$F\in\nabla_\nu(\mathbb Z)$, proving the first equality.

The space of polynomials of the fixed column degrees has a finite
monomial basis. Hence $\nabla_\nu(\mathbb Z)$ is a finitely generated
free abelian group. It spans $V_\nu$, since denominators in a rational
polynomial can be cleared and the defining identities have integral
coefficients. Its subgroup $\Delta_\nu(\mathbb Z)$ is also finitely
generated. It spans $V_\nu$ as well: otherwise nondegeneracy of $B_\nu$
would give a nonzero rational vector $F$ perpendicular to its rational
span. Then every rational multiple of $F$ would belong to
$\Delta_\nu(\mathbb Z)^{\vee,B}=\nabla_\nu(\mathbb Z)$, which is
impossible for a nonzero polynomial with integral coefficients.

Every full lattice in a finite-dimensional rational space with a
nondegenerate symmetric rational bilinear form equals its double dual.
Indeed, take a basis of the lattice; in these coordinates its dual is
the inverse Gram matrix applied to $\mathbb Z^q$, and dualizing once
more returns $\mathbb Z^q$ (the form here is symmetric).
Dualizing the first equality proves the second one. The same coordinate
calculation proves perfection. If $n=0$, both lattices are
$\mathbb Z$, $H_\nu=C_\nu=1$, and these assertions also hold directly.
\end{proof}

\begin{remark}[Saturation has two different roles]\label{coef:saturation}
The lattice $\nabla_\nu(\mathbb Z)$ is saturated in the group of
integral polynomials of column degree $\nu$: if a nonzero integer
multiple of an integral polynomial satisfies the defining identities,
the polynomial itself satisfies them. The smaller lattice
$\Delta_\nu(\mathbb Z)$ need not be saturated in
$\nabla_\nu(\mathbb Z)$.
For binary quadratics, $H=x^2$, $C=2$,
\[
 \nabla=\mathbb Zx^2\oplus\mathbb Zxy\oplus\mathbb Zy^2,
 \qquad
 \Delta=\mathbb Zx^2\oplus2\mathbb Zxy\oplus\mathbb Zy^2.
\]
Replacing $\Delta$ by its saturation would destroy the asserted
integral pairing, since $B(xy,xy)=1/2$.
Also, the class of $2xy$ in $\Delta/2\Delta$ is nonzero although its
image in $\nabla/2\nabla$ is zero. Reducing the inclusion of the
lattices modulo two is not the same as identifying their reductions.
The perfect pairing between the two abstract integral lattices remains
valid at this prime. This explains why a small residue characteristic
does not invalidate Theorem~\ref{coef:integral-duality}, and why the
torsion in $\nabla/\Delta$ must not be discarded.
\end{remark}

\subsection{Number fields, determinant characters and tensor products}

Let $K$ be a number field with ring of integers $\mathcal O$.
The defining conditions for $\nabla_\nu$ are finitely many homogeneous
linear equations with integer coefficients in the finite monomial
basis of column degree $\nu$. Passing from $\mathbb Z$ to $\mathcal O$
or $K$ commutes with their common kernel. One elementary proof puts
the integer matrix of these equations in Smith normal form: each
nonzero diagonal integer acts injectively on either characteristic-zero
domain, so the kernel is spanned by the same remaining basis columns.
Consequently
\[
 \nabla_\nu(\mathcal O)=\mathcal O\otimes_{\mathbb Z}
                 \nabla_\nu(\mathbb Z),\qquad
 \Delta_\nu(\mathcal O)=\mathcal O\otimes_{\mathbb Z}
                 \Delta_\nu(\mathbb Z).
\]
The second identity also follows immediately from the integral formulas
for the divided operators. Extending the two dual integral bases from
Theorem~\ref{coef:integral-duality} proves the same perfect pairing over
$\mathcal O$, with duality taken inside $\nabla_\nu(K)$.

The actual coefficient weights are dominant integral weights
$\lambda_1\ge\cdots\ge\lambda_n$; the $\lambda_i$ may be negative.
For $n>0$ put
\[
 a=\lambda_n,\qquad \nu_i=\lambda_i-a.
\]
Thus $\nu$ is a partition with last part zero. For $n=0$ put $a=0$.
On the square polynomial space define the representation
\begin{equation}\label{coef:representation}
 \rho_\lambda(g)F(X)=\det(g)^a F(g^tX),\qquad g\in\mathrm{GL}_n(K).
\end{equation}
It is a representation, since substitutions reverse matrix
multiplication and the transpose reverses it again. Negative powers
of $\det g$ are legitimate because $g$ is invertible.
Lemma~\ref{coef:adjoint} gives
\begin{equation}\label{coef:contravariant}
 B_\nu(\rho_\lambda(g)F,G)
       =B_\nu(F,\rho_\lambda(g^t)G),\qquad
 B_\nu(\rho_\lambda(g^{-t})F,\rho_\lambda(g)G)=B_\nu(F,G).
\end{equation}

Both standard lattices are preserved by
$\mathrm{GL}_n(\mathcal O)$ and by
$\mathrm{GL}_n(\mathcal O_{\mathfrak p})$ after localization.
For $\nabla$ this follows directly from integral substitution and the
fact that a positive or negative power of a unit determinant is a unit.
For $\Delta=\nabla^{\vee,B}$, it follows from
\eqref{coef:contravariant} and the stability of $\nabla$ under $g^t$;
apply the argument also to $g^{-1}$ to get equality.

For a finite index set $\mathcal T$, take a weight $\lambda_\tau$
for each $\tau$ and use disjoint variable blocks. Let
\[
 V=\bigotimes_{\tau\in\mathcal T}\nabla_{\nu_\tau}(K),\qquad
 N_0=\bigotimes_\tau\nabla_{\nu_\tau}(\mathcal O),\qquad
 D_0=\bigotimes_\tau\Delta_{\nu_\tau}(\mathcal O).
\]
The first tensor product is over $K$ and the latter two are over
$\mathcal O$. They are identified with products of polynomials in
the disjoint blocks. Monomial factorials factor over the blocks, so
the tensor-product pairing is exactly the Fock pairing divided by
\[
                             C=\prod_\tau C_{\nu_\tau}.
\]
Tensoring the dual bases proves
\begin{equation}\label{coef:standard-tensor-duality}
                  N_0^{\vee,B}=D_0,\qquad D_0^{\vee,B}=N_0.
\end{equation}
The representation $\rho(g)$ is the tensor product of
\eqref{coef:representation}. Formula~\eqref{coef:contravariant}
continues to hold componentwise. For an empty index set the space is
$K$, both lattices are $\mathcal O$, and $C=1$.

The integral bases constructed above give rational polynomial bases of
$V$. The matrix of $B$ in such a basis has rational entries. If all
entries of the matrices $g_\tau$ lie in a subfield $K_0\subset K$,
the matrix of $\rho(g)$ in the same basis also has entries in $K_0$:
the polynomial substitution and determinant powers have coefficients
in $K_0$, and expressing their outputs in a rational basis uses only
linear equations with rational coefficients. This gives the required
field of definition without making a special tableau-basis choice.

\subsection{Transport by the actual local bases}

Let $M_\tau\subset K^n$ be full $\mathcal O$-lattices. For every
nonzero prime $\mathfrak p$, choose matrices
$B_{\tau,\mathfrak p}\in\mathrm{GL}_n(K)$ whose columns form bases
of $(M_\tau)_{\mathfrak p}$ over $\mathcal O_{\mathfrak p}$.
Put $B_{\mathfrak p}=(B_{\tau,\mathfrak p})_\tau$ and define
\begin{equation}\label{coef:transport}
 N(M)=\bigcap_{\mathfrak p}\rho(B_{\mathfrak p})(N_0)_{\mathfrak p},
 \qquad
 D(M)=\bigcap_{\mathfrak p}\rho(B_{\mathfrak p})(D_0)_{\mathfrak p}
                         \quad\text{inside }V.
\end{equation}

\begin{lemma}[Well-defined transported lattices]\label{coef:transport-lattices}
The intersections in \eqref{coef:transport} are independent of the
chosen local bases. They are full lattices, their localizations are
the corresponding displayed local modules, and $D(M)\subset N(M)$.
For a tuple $g=(g_\tau)$ of invertible matrices,
\begin{equation}\label{coef:equivariance}
                  N(gM)=\rho(g)N(M),\qquad D(gM)=\rho(g)D(M).
\end{equation}
If $M_\tau^\vee$ denotes the ordinary bilinear dual lattice in $K^n$,
then
\begin{equation}\label{coef:transport-duality}
                     N(M)^{\vee,B}=D(M^\vee),
              \qquad D(M)^{\vee,B}=N(M^\vee).
\end{equation}
\end{lemma}
\begin{proof}
Two bases differ on the right by an element of
$\mathrm{GL}_n(\mathcal O_{\mathfrak p})$, which preserves the two
standard localized lattices. This proves independence.
Outside a finite set of primes the base lattices are
$\mathcal O_{\mathfrak p}^n$, so the identity matrices can be chosen.
The local modules in \eqref{coef:transport} therefore equal the
standard ones outside a finite set.

Here is the elementary patching argument. Given finitely many
exceptional local lattices and a reference lattice $L_0$, choose a
nonzero $b\in\mathcal O$ such that each lies between
$b(L_0)_{\mathfrak p}$ and $b^{-1}(L_0)_{\mathfrak p}$.
The finite module $b^{-1}L_0/bL_0$ is the direct sum of its primary
parts. In each part take the submodule specified by the corresponding
local lattice; at the additional primes dividing $b$ take the
submodule specified by $L_0$. Its inverse image in $b^{-1}L_0$ is a
full lattice with exactly the prescribed localizations. A lattice
equals the intersection of its localizations: in any finite generating
set, a failure of membership is detected in a primary component of
a finite quotient after clearing denominators. This proves both the
existence assertion and its equality with the intersections above.
The inclusion follows from $D_0\subset N_0$.
The matrices $g_\tau B_{\tau,\mathfrak p}$ are bases of
$(g_\tau M_\tau)_{\mathfrak p}$, proving
\eqref{coef:equivariance}.

Ordinary duality of a full lattice commutes with localization. For
example, an element pairing locally integrally with a finite generating
set can be multiplied by a single element outside $\mathfrak p$ so
that it pairs globally integrally with all generators. The bases of
$(M_\tau^\vee)_{\mathfrak p}$ are $B_{\tau,\mathfrak p}^{-t}$.
Equations~\eqref{coef:contravariant} and
\eqref{coef:standard-tensor-duality} therefore give
\[
 \bigl(\rho(B_{\mathfrak p})(N_0)_{\mathfrak p}\bigr)^{\vee,B}
     =\rho(B_{\mathfrak p}^{-t})(D_0)_{\mathfrak p}.
\]
Intersecting over all primes proves the first equality in
\eqref{coef:transport-duality}; exchange $N_0,D_0$ for the second.
\end{proof}

The definition in \eqref{coef:transport} is exactly the intersection
over all finite primes and all admissible local bases used in the
coefficient construction. The independence just proved is why those
extra quantifiers do not shrink the lattice further.

\subsection{The normalized Hermitian form and unimodular inputs}

Suppose $K$ has an involution $c\mapsto\bar c$.
It preserves $\mathcal O$ and acts coefficientwise on polynomials.
For each $\tau$ let $G_\tau\in\mathrm{GL}_n(K)$ be Hermitian,
$\bar G_\tau^{\,t}=G_\tau$. On $V$ define
\begin{equation}\label{coef:hermitian}
 h_G(F,F')=B(\bar F,\rho(G)F').
\end{equation}
Written in terms of polynomials, this is exactly
\[
 h_G(F,F')=
 \frac{\prod_\tau\det(G_\tau)^{a_\tau}}
      {\prod_\tau C_{\nu_\tau}}
 [\bar F,\ F'((G_\tau^tX_\tau)_\tau)].
\]
Thus both the factorial normalization and the determinant characters
are retained. Formula~\eqref{coef:contravariant} shows that $h_G$ is
Hermitian and nondegenerate. At a complex embedding where every
$G_\tau$ becomes positive definite, its numerical Gram matrix in a
rational polynomial basis is positive definite as well. Indeed the
ordinary complex Fock form is positive definite, the adjoint of
$\rho(g)$ is $\rho(g^*)$, and a factorization $G=SS^*$ gives
$\rho(G)=\rho(S)\rho(S)^*$.
The same adjoint identity gives, componentwise,
\[
 h_G(\rho(g)F,\rho(g)F')=h_{g^*Gg}(F,F').
\]
Thus a base isometry $g^*Gg=G$ acts by isometries on the coefficient
space. If it also preserves the input lattices, it preserves both
coefficient lattices by \eqref{coef:equivariance}. These are the
isometry and stability properties used for their finite stabilizers.

For a base lattice put
\[
 M_\tau^{\#,G_\tau}
   =\{y:\bar x^{\,t}G_\tau y\in\mathcal O\text{ for all }x\in M_\tau\}
   =G_\tau^{-1}(\bar M_\tau)^\vee.
\]
The symbol $\#$ on coefficient lattices below refers to $h_G$.

\begin{theorem}[The coefficient duality actually needed]\label{coef:hermitian-duality}
With all the preceding definitions,
\begin{equation}\label{coef:hermitian-duality-eq}
             N(M)^\#=D(M^{\#,G}),\qquad
             D(M)^\#=N(M^{\#,G}).
\end{equation}
In particular, if every input lattice is unimodular for its base
Hermitian form, then
\begin{equation}\label{coef:unimodular}
                         N(M)^\#=D(M)\subset N(M),\qquad
                         D(M)^\#=N(M).
\end{equation}
Consequently these coefficient lattices satisfy the initial dual
containments in the finite-group argument with $s=1$.
\end{theorem}
\begin{proof}
All standard lattice generators have integer coefficients. Conjugation
therefore gives $\overline{N(M)}=N(\bar M)$ and
$\overline{D(M)}=D(\bar M)$: conjugating a local basis at
$\overline{\mathfrak p}$ gives the local basis at $\mathfrak p$.
For $y\in V$, the defining condition for $N(M)^\#$ is
\[
 B(\overline{N(M)},\rho(G)y)\subset\mathcal O.
\]
By symmetry of $B$ and \eqref{coef:transport-duality}, this is
equivalent to $\rho(G)y\in D((\bar M)^\vee)$.
Equivariance \eqref{coef:equivariance} makes the resulting lattice
\[
 \rho(G^{-1})D((\bar M)^\vee)
                   =D(G^{-1}(\bar M)^\vee)=D(M^{\#,G}).
\]
The other equality follows by exchanging the two lattices.
The unimodular conclusion is immediate, including its inclusion,
which was proved in Lemma~\ref{coef:transport-lattices}.
\end{proof}

\begin{lemma}[Extension of a unimodular input]\label{coef:extension}
Let $E\hookrightarrow K$ be an embedding of number fields compatible
with their involutions. For a nondegenerate Hermitian $E$-space and a full
$\mathcal O_E$-lattice $L$, extension of scalars satisfies
\[
 (\mathcal O_K L)^\#=\mathcal O_K L^\#,
\]
where the form on the left is the scalar extension of the original
$E$-valued form. In particular a unimodular lattice remains unimodular,
even at ramified primes.
\end{lemma}
\begin{proof}
The module $L$ is finite projective over the Dedekind domain
$\mathcal O_E$. The torsion-free module $\mathcal O_K$ is flat over
$\mathcal O_E$. Duals of finite projective modules commute with scalar
extension: this is immediate for a finite free module, and follows
for a projective module by writing it as a direct summand of one.
Compatibility of the involutions also identifies the extension of
$\bar L$ with the conjugate of its extension. Use the formula
$L^\#=G^{-1}(\bar L)^\vee$ in any coordinates over $E$ and extend
scalars. The pairing is extended as a field-valued pairing; no trace
map is inserted, so there is no different-ideal factor.
\end{proof}

The first coefficient construction takes these scalar extensions of
the original unimodular base lattices. The second uses their ordinary
duals. The latter are unimodular for the inverse
transpose Gram form, because direct calculation gives
\[
 (M^\vee)^{\#,(G^{-1})^t}=(M^{\#,G})^\vee.
\]
Indeed, both sides equal $G^t\bar M$, using
$(\bar M^\vee)^\vee=\bar M$.
Thus Theorem~\ref{coef:hermitian-duality} supplies the exact duality
for both constructions once their stated base-lattice unimodularity is
available. The construction and uniform estimates for those underlying
arithmetic lattices remain separate requirements.

\begin{example}[A determinant character of negative weight]
For $n=1$ and weight $\lambda=-1$, the polynomial partition is zero
and $V=K$, but $\rho(g)=g^{-1}$.
Take $K=\mathbb Q(i)$, $M=\mathcal O$ and $G=2$.
Then $M^{\#,G}=\tfrac12\mathcal O$,
$N(M^{\#,G})=2\mathcal O$, and
$h_G(x,y)=\bar x y/2$ has dual lattice $2\mathcal O$, as required.
If the determinant factor in \eqref{coef:hermitian} were omitted,
the coefficient dual would instead be $\mathcal O$.
Negative weights therefore cannot be replaced by their nonnegative
polynomial parts without retaining the determinant character.
More generally, for any fractional ideal $I$ and integral rank-one
weight $a$, local transport gives $N(I)=D(I)=I^a$, including negative
$a$ and nonprincipal $I$. The Hermitian dual for Gram entry $g$ is
$g^{-a}\bar I^{-a}$, equal to $D(g^{-1}\bar I^{-1})$ as the theorem
asserts. This also checks the transport convention when no global
basis of the input lattice exists.
\end{example}

\subsection{Checks and the resulting scope of the reconstruction}

An accompanying exact computation constructs the two lattices
independently for twelve small cases. The dual-Weyl lattice is the
saturated integer kernel of the full column-invariance equations;
the Weyl lattice is the lattice generated from $H_\nu$ by closure
under all divided row operators. Smith and Hermite normal forms are
used over the integers. In every case the normalized matrix pairing
the two lattice bases is integral and has determinant $\pm1$.
For example, for square matrices the partitions $(2,1,0)$ and
$(3,1,0)$ have ranks $8$ and $15$, normalizing constants $3$ and $8$,
and indices $[\nabla:\Delta]=3$ and $4096$, respectively.
For four rows, two columns and partition $(2,1)$, the rank is $20$
and the index is $81$. These finite checks support the stated
conventions; the proof of the general result is the preceding argument.

The coefficient-lattice duality, its determinant twists, tensor
products, local transport and scalar extension are now proved in
ordinary mathematics with the actual definitions used by the divisor
construction. The proof uses characteristic zero, $n\le r$, a
nonincreasing nonnegative polynomial partition, invertible base Gram
matrices and full input lattices. The original integral weights may
be arbitrary dominant integers. No assumption that the residue
characteristic exceeds the degree is made. No counterexample to the
actual duality has been found; the explicit examples above instead
identify changes of normalization or saturation that would make it
false.



% Exact Fock normalizer and local equality of Weyl and dual-Weyl lattices.
% Labels in this fragment have the prefix fh-.
